\section{A.1. What the surgery maps can and cannot mean}
\subsection{The target in the manuscript}
On cote 67, scans 074--078, the normalization of the stable curve carries at every node a nonzero element of the tensor product of the two branch tangent lines. Scan 078 also uses its inverse as a perfect pairing. We denote the tensor by $q_e\in(T_{a_e}Y\otimes_{\C}T_{b_e}Y)^\times$. No bound on $q_e$ is imposed there. We first use this literal target. All curves in this section are complex, all marked points are distinct, and $2g-2+n>0$.

Two independent definition choices affect the arrow: which cuts are
allowed, and which bordered data the target retains. For actual
arbitrary smooth embedded cuts, the two canonical normalized
arclength structures need not agree;
Proposition~\ref{A1-nontoric-cut} gives an analytic counterexample
to the asserted cap antiisomorphism. Geodesic representatives of
cut isotopy classes give toric identifications with equal paired
hyperbolic lengths. Requiring only toric compatibility is a broader
choice: it allows unequal lengths. Theorem~\ref{A1-toric-domain}
proves that this broader map has the same tensor and shadow image
but an additional real parameter for every seam, and
Proposition~\ref{A1-toric-noninjective} gives an explicit
four-pointed counterexample to injectivity.

In the bordered case, requiring the retained external discs to be
admissible after the node branches are left as punctures fails by
Proposition~\ref{A1-bordered-landing}. One can instead require
joint admissibility of all caps, or record their external shadows
and reconstruct the full systems jointly. The equal-length model
is described by Theorems~\ref{A1-graph-range} and
\ref{A1-bordered-shadow-range}, with the full-disc criterion in
Corollary~\ref{A1-bordered-shape-image}. For the broader toric
model, Corollary~\ref{A1-toric-full-discs} gives the corresponding
criterion without asserting that full external shapes have the
same fibers as their shadows. The tensor image is proper whenever
there is an edge. The cap-size obstruction applies to every
interpretation in which the tensor is produced by the stated
toric cap construction.

\subsection{The choice between embedded cuts and geodesic cuts}
\label{sec:A1-cut-choice}

\begin{proposition}[Refuted: arbitrary analytic cuts canonically give toric cap identifications]
\label{A1-nontoric-cut}
If the cut circles in 67/074 are interpreted as arbitrary embedded
smooth, or even analytic, circles, cutting and canonical capping
does not always give the asserted antiisomorphisms of pointed caps.
This failure already occurs on a four-punctured sphere with one
separating circle. The geodesic-cut model used in
Theorems~\ref{A1-graph-range} and~\ref{A1-bordered-shadow-range}
therefore requires geodesic representatives, or toric compatibility
together with equality of the paired hyperbolic boundary lengths.
Smoothness alone does not supply these conditions. Toric compatibility
without length equality defines the larger, noninjective model of
Theorem~\ref{A1-toric-domain}.
\end{proposition}
\begin{proof}
Take two hyperbolic pants, each with two cusp ends and one
geodesic boundary of length $2\pi$. Let $w_1,w_2$ be their canonical
cap coordinates, so that $w_j=e^{i\theta_j}$ parametrizes the
boundary by normalized arclength. Glue these boundaries by
\[
 w_2=h(w_1),\qquad
 h(e^{i\theta})=
       \exp\bigl(i(-\theta+\tfrac12\sin\theta)\bigr).
\]
The angular derivative is $-1+\tfrac12\cos\theta$, everywhere
between $-3/2$ and $-1/2$. Thus $h$ is an orientation-reversing
analytic diffeomorphism of the circle. It extends to the holomorphic
annulus map
\[
 H(w)=w^{-1}\exp\bigl((w-w^{-1})/4\bigr).
\]
Its derivative does not vanish on the circle. Since its restriction
to that circle is injective, compactness and local injectivity imply
that it is injective on a sufficiently small whole annulus. Gluing by
this collar extension gives a Riemann surface with an analytic
seam, topologically a four-punctured sphere. Equivalently this
is the analytic special case of Theorem~\ref{A3-welding}.

Cutting along that seam recovers the two original conformal
pants. Their unique canonical hyperbolic metrics, and hence their
canonical cap coordinates, are exactly those used above. An
antiisomorphism of centered unit discs has the form
$w\mapsto c\overline w$, $|c|=1$, by the Schwarz lemma. Its boundary
angular derivative is constantly $-1$. The identification $h$
has derivatives $-1/2$ at $\theta=0$ and $-3/2$ at $\theta=\pi$,
so it is not the boundary map of any such pointed-disc
antiisomorphism. The tensor claimed in 67/075 cannot therefore
be extracted by that stated construction from this analytic cut.
\end{proof}

If the cut datum means an essential isotopy class instead, the
canonical hyperbolic metric of the entire surface selects its
unique geodesic representative. Simultaneously disjoint classes
have disjoint geodesic representatives. Cutting along those
representatives restricts the same hyperbolic metric to both
banks, so the gluing is an arclength isometry. This is the
geodesic interpretation for which the exact range theorems
above apply. It changes the embedded analytic seam of the
counterexample and must be named as part of the definition.


\begin{proposition}[Known with derivation: the dimension comparison]
\label{A1-dimensions}
We use orientation-preserving markings that preserve the peripheral
classes and labelled boundary circles; boundary points may move during
isotopy. The positive-boundary Teichm\"uller space of stable type $(g,n)$
has real dimension $6g-6+3n$. For a connected stable graph of total genus
$g$ with $n$ legs, its normalized component moduli together with a
nonzero complex tangent tensor at every edge have complex dimension
$3g-3+n$. Adding $n$ external positive cap parameters gives real
dimension $6g-6+3n$. These counts verify 67/077 and 67/016, but do not
establish existence or surjectivity of the proposed arrows.
\end{proposition}
\begin{proof}
The bordered-space convention and real dimension are also stated in
Bers's original account, \cite[\S2(F), p.~139; \S3(F), p.~141;
\S4(E), p.~143]{ABers1981}: boundary points are not marked pointwise,
and the bordered extremal problem is obtained by doubling. The
Fenchel--Nielsen coordinates recalled in Mondello's subsection
``Teichm\"uller space'' \cite{AMondello2011} consist of one length and
one twist for each of $3g-3+n$ interior pants curves, and one positive
length for each of the $n$ labelled boundary circles. Their number is
$2(3g-3+n)+n$. For the graph let $V,E$ be its vertices and edges, and
let $g_v,n_v$ be the component genus and number of flags. Then
$\sum_v n_v=2|E|+n$ and
$g=\sum_vg_v+|E|-|V|+1$. The complex dimension is consequently
\[
 \sum_{v\in V}(3g_v-3+n_v)+|E|=3g-3+n.
\]
The $n$ real cap parameters are supplied by
Theorem~\ref{A4-grafting}, so the final count follows. Finite graph
automorphisms do not change dimensions. The formula includes the
zero-dimensional unbordered $(0,3)$ case and the one-dimensional
unbordered $(1,1)$ case.

The marking qualification is substantive. On the three-punctured
sphere, the free-group automorphism $a\mapsto ab$, $b\mapsto b$ is a
homotopy equivalence, with inverse $a\mapsto ab^{-1}$, $b\mapsto b$.
It sends the peripheral class $(ab)^{-1}$ to $(ab^2)^{-1}$, whose
abelianization $-(1,2)$ is not that of any peripheral class. Thus an
arbitrary homotopy equivalence is not the marking used in these moduli
spaces. This verifies the input correction credited to Michel Hua's
project without taking it on authority.
\end{proof}

\begin{lemma}[Known, with independent proof: the cap bound]\label{A1-cap-bound}
Let $f:\D\longrightarrow\C\setminus\{1\}$ be injective and holomorphic, with $f(0)=0$. Then $|f'(0)|\leq4$. Thus a pointed disc in $\mathbf P^1$ centred at $0$ and avoiding $1,\infty$ has conformal radius at most $4$ in the coordinate $z$.
\end{lemma}
\begin{proof}
We include the classical quarter-theorem argument. The same sharp
bound and slit equality case are explicitly stated and proved in the
opened original \cite[\S2, pp.~173--174]{ALavrentiev1934}.
For historical precision, Koebe's opened 1907 original gives a positive
universal covering radius in its second auxiliary theorem and its
alternative form, pp.~203--205 \cite{AKoebe1907}. Those pages do not
state the sharp constant $1/4$; the sharp estimate used here is proved
in the following argument and in the cited Lavrentiev passage.
If $F(z)=z+a_2z^2+\cdots$ is univalent in $\D$, the exterior odd function
\[
G(w)=\frac{1}{\sqrt{F(w^{-2})}}=w-\frac{a_2}{2w}+\cdots\qquad(|w|>1)
\]
is univalent. The square root is chosen by its expansion at infinity; injectivity follows because equality of squares forces $w_1=\pm w_2$, while oddness excludes the minus sign. For an exterior univalent function $w+b_0+\sum_{k\geq1}b_kw^{-k}$, the area enclosed by the image of $|w|=r>1$ is
\[
\pi\left(r^2-\sum_{k\geq1}k|b_k|^2r^{-2k}\right)\geq0.
\]
The equality follows by integrating $\overline G\,dG/(2i)$, with uniformly convergent Laurent series on that circle. Letting $r\downarrow1$ gives the area theorem. Applied to $G$ it gives $|a_2|\leq2$. If $a\ne0$ is omitted by $F$, the function $F/(1-F/a)$ is again normalized and univalent, and has second coefficient $a_2+1/a$. Hence $|1/a|\leq4$. Apply this to $F=f/f'(0)$ and the omitted value $a=1/f'(0)$.
\end{proof}
\begin{lemma}[Proved: the tensor produced by a toric seam]\label{A1-tensor}
Let $f:\overline\D\to D_a\subset Y$ and $h:\overline\D\to D_b\subset Y$ parametrize two pointed analytic discs, with centres $a,b$. An orientation-reversing toric identification of their boundaries has the form $f(\zeta)\sim h(c/\zeta)$, $|\zeta|=1$, where $|c|=1$. Its tangent tensor is
\[
q=c\,df_0(\partial_\zeta)\otimes dh_0(\partial_\eta).
\]
In any chosen branch coordinates its modulus is the product of the two conformal radii.
\end{lemma}
\begin{proof}
The annular continuation of the boundary identification is $\zeta\eta=c$. Applying the derivatives to the two tangent factors gives the displayed tensor. Replacing $f(\zeta)$ by $f(\alpha\zeta)$ and $h(\eta)$ by $h(\beta\eta)$ changes $c$ to $c/(\alpha\beta)$ and the derivative factors to $\alpha df_0,\beta dh_0$, so the tensor is independent of the two rotations. Taking absolute values proves the last assertion.
\end{proof}

\begin{theorem}[Refuted: the global isomorphisms in A.1 and L2]\label{A1-global-refutation}
With the unrestricted nonzero-tensor target of cote 67/075--078,
the proposed global isomorphisms are false. For the geodesic-cut
interpretation the unbordered map is not surjective on complex objects,
already for $(g,n)=(0,4)$ and then for every stable type with
$3g-3+n>0$. The same cap-size obstruction rules out a preimage for
specified bordered target objects whenever the tensor is obtained
from the actual caps. This is an obstruction to coarse surjectivity
as well as to equivalence of stacks or groupoids. Under the external-only
admissibility reading the bordered arrow also has the separate failure
of being undefined on some objects, proved in
Proposition~\ref{A1-bordered-landing}; it must not be described simply
as a globally defined map with a proper image. For arbitrary embedded
cuts, Proposition~\ref{A1-nontoric-cut} supplies the earlier failure
of the tensor construction itself.
\end{theorem}
\begin{proof}
For $(0,4)$ take two copies of $\mathbf P^1$, identify their points $0$ to form a node, and use $1,\infty$ on each as the four labelled markings. This is a connected stable curve of arithmetic genus zero with four markings. Put
\[
q=17\,\partial_z|_0\otimes\partial_w|_0.
\]
It is a valid nonzero nodal tangent tensor. Any preimage under the asserted surgery would give, on each normalized sphere, an embedded pointed cap about $0$ avoiding $1$ and $\infty$. Lemmas \ref{A1-cap-bound} and \ref{A1-tensor} would imply $|q|\leq4\cdot4=16$, contrary to its definition. The markings fix the chosen normalization, so this contradiction is invariant under isomorphism.

For every stable $(g,n)$ of positive complexity, a pants decomposition gives a connected stable trivalent graph with $n$ legs, first Betti number $g$, and $3g-3+n>0$ edges. Replace each vertex by a sphere with its three flags at $0,1,\infty$. At each flag use a local projective coordinate sending that flag to $0$ and the other two flags to $1,\infty$. The ambiguity of interchanging the latter two is $z\mapsto z/(z-1)$, whose derivative at zero has modulus one. Thus each flag has a well-defined tangent norm, preserved by all automorphisms of the graph curve. Choose a nonzero tensor of norm $17$ at every edge. A preimage would require all these norms to be at most $16$. This is impossible, including for a self-node, such as the one-node genus-one curve with one marking.

For the external-only bordered target also choose admissible external
discs. These exist: at each trivalent vertex take a generalized
hyperbolic pant with cusp ends at the incident nodes and arbitrary
positive lengths at the incident external legs; canonically cap just
its actual boundaries. The resulting external discs are admissible
on the sphere with its node branches left as marked punctures.
The node tensors in that stated target are independent extra data,
so choose their norms to be $17$ as above. Any hypothetical preimage
would still have embedded node caps avoiding the other special
points. This gives the same contradiction, even if the source is
restricted to the objects whose external output happens to satisfy
external-only admissibility. In the repaired joint target the
support inequalities explicitly exclude these tensors.
\end{proof}

\begin{proposition}[Proved: the zero-complexity edge case]\label{A1-zero}
For $(g,n)=(0,3)$ there is no essential stable cut. The unbordered groupoid with labelled markings is a point. The bordered capping functor is an equivalence with pointed spheres equipped with an admissible three-disc system, when admissible has its definition from 67/065--066 and arrows retain those discs.
\end{proposition}
\begin{proof}
The sphere with three ordered points has a unique projective normalization and no nontrivial automorphism fixing all three. For the bordered assertion, remove the three open discs. Admissibility says exactly that canonical arclength capping reconstructs the same pointed discs. Conversely, removing the caps from a canonically capped pant recovers the pant. Conformal isomorphisms restrict and extend uniquely across the analytic seams. These are inverse functors, not an identification of arbitrary radii with discs.
\end{proof}

\begin{proposition}[Proved: the corrected surgery equivalence]\label{A1-retained-discs}
Fix a stable graph. Enlarge the target by retaining admissible pointed caps at every branch of every node, as well as at the external markings, and a toric boundary identification at every node. Require the two hyperbolic boundary lengths at a node to be equal. Then geodesic cutting and canonical capping give an equivalence of groupoids with hyperbolic bordered surfaces cut along that graph. The nodal tensor alone is a forgetful invariant of this equivalence.
\end{proposition}
\begin{proof}
Remove all open caps. Each component has its unique hyperbolic metric with geodesic boundary. By admissibility its boundary circle actions are the given ones. A toric orientation-reversing identification between circles of equal circumference is an isometry. In signed normal distance $r$ the collar metric is $dr^2+\cosh^2(r)\,ds^2$; therefore the two collars glue to a smooth curvature-minus-one metric with the seam a geodesic. Its conformal structure gives the required welded surface. Cutting along the indicated seams recovers the original pieces and caps. Both operations carry isomorphisms to isomorphisms and are mutually inverse. If there are punctures, they retain their complete cusp neighborhoods throughout. The uniqueness of the conformal structure can also be read from analytic reflection across these seams.
\end{proof}

\subsection{The exact unbordered range on each stable-graph stratum}
\label{sec:A1-range}

This subsection concerns the unbordered construction: the external
marked points remain cusps. Let $\Gamma$ be a stable graph with edge
set $E$, and fix its marked normalization $(Y_v)_{v\in V}$.
At a vertex, $n_v$ counts all node branches and external markings;
assume $2g_v-2+n_v>0$. Write $I_v$ for its node branches, and let
\[
 A_v:\R^E\longrightarrow\R^{I_v},\qquad
 (A_vs)_a=s_{e(a)}.
\]
A loop edge has two distinct branches at the same vertex and therefore
occurs twice among these coordinates. For each edge the two geodesic
boundary lengths are required to be equal, with common value
$2\pi\sqrt{s_e}$.

Let $S_v$ and $K_v$ be the potential and support function of
Theorems~\ref{A2b-energy}--\ref{A2b-image}, applied to $Y_v$ with
active set $I_v$ and the other markings as cusps. A vertex with no
active branches contributes zero. Define
\begin{equation}\label{A1-graph-potential}
 \Psi_\Gamma(s)=\sum_v S_v(A_vs),\qquad
 K_\Gamma(p)=\sum_v K_v(A_vp).
\end{equation}
All logarithmic norms below use arbitrary fixed tangent frames at the
branches; changing the frames translates their image in $\R^E$.

\begin{theorem}[Proved: exact range, tangent map, and fundamental group]
\label{A1-graph-range}
For the fixed normalization, the unbordered canonical surgery map,
with actual toric seam identifications as its phase variables, is a
real-analytic diffeomorphism onto
\[
 \mathcal U_\Gamma(Y)=
 \left\{(q_e)\in\prod_{e\in E}(T_{a_e}Y\otimes T_{b_e}Y)^\times:
   \sum_e p_e\log|q_e|<K_\Gamma(p)
   \quad\text{for all }0\ne p\in[0,\infty)^E\right\}.
\]
Its logarithmic radial image is open, convex, and downward closed.
If $E\ne\varnothing$, this is a proper subset of the unrestricted
nonzero-tensor fiber. The tangent map of surgery is an isomorphism.

As the normalization varies on its fixed graph stratum, these images
form an open real-analytic subbundle of the nonzero-tensor bundle.
Its inclusion is an equivariant homotopy equivalence, and consequently
induces an isomorphism on fundamental groups, including orbifold
fundamental groups when finite automorphisms are retained. These
conclusions allow loop edges as well as edges between different
components.
\end{theorem}
\begin{proof}
At the branch $a$ of an edge $e$, let $r_a(s)$ denote the canonical
cap radius obtained on $Y_{v(a)}$ with parameters $A_{v(a)}s$.
The tensor computation in Lemma~\ref{A1-tensor} gives
\[
 \log|q_e|=\log r_{a_e}(s)+\log r_{b_e}(s)
          =\frac{\partial\Psi_\Gamma}{\partial s_e}(s).
\]
This chain rule counts both branches of a loop separately. Its Hessian
is
\[
 d^2\Psi_\Gamma=\sum_v A_v^T(d^2S_v)A_v.
\]
Each nonempty term is positive definite on the branch variables. If
$\dot s\ne0$, choose an edge with $\dot s_e\ne0$; at least one
$A_v\dot s$ is nonzero. Hence the displayed Hessian is positive
definite on $\R^E$. The logarithmic tensor map is therefore locally
invertible and globally injective by integration along straight
segments, exactly as in Theorem~\ref{A2b-energy}.

For completeness the same argument gives the entire stated image,
not just an injective parametrization. Its recession function is
$K_\Gamma$: apply Lemma~\ref{A2b-recession} to each $A_vs_0+tA_vp$
and sum. The supports are finite and positively homogeneous, with
$K_v(0)=0$. Every actual radius vector satisfies all the strict
inequalities, because at least one nonzero $A_vp$ gives a strict
component inequality.

Conversely, let $x\in\R^E$ satisfy those strict inequalities.
The convex function $\Psi_\Gamma(s)-x\cdot s$ has bounded sublevels:
an escaping sequence, normalized in direction, would give a nonzero
$p\geq0$ with recession slope $K_\Gamma(p)-p\cdot x\leq0$, by
the segment argument in the proof of Theorem~\ref{A2b-image}.
Its lower semicontinuous extension to the closed orthant attains its
minimum. A finite boundary minimum is impossible. Along an inward
diagonal, every vanishing edge variable produces two cap terms tending
to $-\infty$ by the uniform collar bound, while all other cap terms
have uniform finite upper bounds. The inward derivative thus tends
to $-\infty$. The minimum is interior, unique, and has gradient $x$.
This proves the asserted logarithmic range and its analytic inverse.

Normalize the cap maps to have positive derivatives in the chosen
tangent frames. The gluing identification is then
$w_{a_e}w_{b_e}=c_e$, $|c_e|=1$, and the argument of $q_e$ is
precisely that of $c_e$. Thus all phases occur independently and
uniquely. Equal lengths make this identification an isometry, so
Proposition~\ref{A1-retained-discs} gives the unique welded hyperbolic
surface with its indicated cut. This proves the real-analytic
diffeomorphism. If the normalization also varies, the map preserves
the normalization coordinate. Its derivative is block triangular
with identity on those coordinates, the positive Hessian on length
variables, and identity on phase variables. It remains invertible;
the local inverses fit together by the fiberwise uniqueness.

To see properness when there is an edge, take $p$ to be its coordinate
unit vector. The corresponding inequality bounds $\log|q_e|$ above
by the finite number $K_\Gamma(p)$. The unrestricted tensor fiber
has no such bound.

Finally set $s_e=1$ for every edge. The canonical caps of these common
lengths determine an intrinsic circle in each nodal tensor line, by
taking the product of their shadow boundary circles. Let $\mathcal C$
be the product of these circles over the normalization stratum.
For any nonzero tensor tuple, let $a_e>0$ be its radial ratio to the
corresponding circle. The homotopy
\[
 H_t((q_e))=(a_e^{-t}q_e)_{e\in E},\qquad0\leq t\leq1,
\]
is a strong deformation retraction of the full nonzero-tensor bundle
onto $\mathcal C$. It also retracts $\mathcal U_\Gamma$ onto
$\mathcal C$, because in logarithmic coordinates it follows a segment
to an interior point of a convex fiber. The construction is intrinsic
and respects graph automorphisms, markings, and relabellings. Both
retractions descend to the relevant orbifold quotients. Their common
retract proves the homotopy and fundamental-group assertions.
For a fixed normalization this reduces to the identity on the
$\Z^E$ fundamental group of the phase torus.
\end{proof}

\begin{corollary}[Proved: the two smallest one-node ranges]
\label{A1-small-ranges}
For the separating one-node $(0,4)$ curve, normalize each rational
component by putting its node at $0$ and its two external markings at
$1,\infty$. In the resulting tensor frame the exact range is
\[
                         0<|q|<16.
\]
For the irreducible one-node $(1,1)$ curve, put its two branches at
$0,1$ and its marking at $\infty$ on its rational normalization.
In the frame $\partial_z|_0\otimes\partial_z|_1$ the exact range is
\[
                         0<|q|<1.
\]
In both cases every phase occurs and the length is uniquely determined
by $|q|$.
\end{corollary}
\begin{proof}
For $(0,4)$ each of the two components has one active cap. Its support
value at weight $1$ is $\log4$ by the Koebe limit in
Corollary~\ref{A2-explicit-limits}. Thus
$K_\Gamma(p)=p\log16$, and Theorem~\ref{A1-graph-range} gives the
first range. Equivalently the two identical one-cap radii increase
from $0$ to $4$ with their common length.

For $(1,1)$ there is one component and a loop edge. Its incidence map
is $s\mapsto(s,s)$. The two-weight Strebel discs at $0,1$ are the
half-planes separated by $\operatorname{Re}z=1/2$, each of conformal
radius $1$. Hence $K_v(p,p)=0$ and the logarithmic range is
$\log|q|<0$. This proves the second assertion. Strict monotonicity
is the one-dimensional positive Hessian of $S_v(s,s)$, and the
phase statement is the seam tensor formula.
\end{proof}

\paragraph{Marking convention and bordered scope.}
The phase in this statement is an actual orientation-reversing toric
identification of the two boundary circles, so it belongs to $S^1$.
If a fully marked Fenchel--Nielsen twist is instead kept as a real
number, its full Dehn-twist translates give the same tensor phase;
one must take that quotient, or lift the tensor phase to the matching
covering space, before asserting the displayed diffeomorphism.

For the bordered target in the manuscript, the external admissible
disc domains are retained as data, not merely their boundary lengths.
Fixing their lengths while varying internal node lengths generally
changes those external domains, by the cap-dependence result in
A.2(a). Therefore Theorem~\ref{A1-graph-range} does not by itself
identify the image in a fiber where the full external domains have
been fixed. Proposition~\ref{A1-bordered-landing} proves the further
obstruction that external-only admissibility need not survive
discarding the node caps. Theorem~\ref{A1-bordered-shadow-range}
instead includes external shadows among the joint gradient variables,
and Corollary~\ref{A1-bordered-shape-image} gives the resulting exact
criterion when the external domains themselves are prescribed.

\subsection{The bordered landing obstruction and its exact repair}
\label{sec:A1-bordered}

There is a separate issue before the size of the tensor range is
considered. Cote 67/065--066 defines admissibility by comparing the
circle actions of the discs with those of the hyperbolic metric on
their complement. The proposed target in 67/075--076 retains an
admissible system only at the external markings. On a nodal curve,
the direct extension of that definition is componentwise on the
normalization, with the branches of the nodes left as punctures.
The following obstruction uses precisely this interpretation.

\begin{proposition}[Refuted: the bordered map lands in the stated target]
\label{A1-bordered-landing}
Canonical cutting and capping need not produce an admissible system
of external discs on the normalization when the node caps are
forgotten and their centers are left as punctures. This already
fails for a four-holed sphere with one separating cut, all five
geodesic lengths---the four external boundaries and the cut---equal
to $2\pi$. Thus, with the componentwise admissibility just specified,
the displayed arrow $MDB\longrightarrow M_{\mathrm{st}}ND$ of
67/076 is not a well-defined map on all objects.
\end{proposition}
\begin{proof}
Take two hyperbolic pants whose three boundary lengths are $2\pi$,
and glue one boundary of each by any orientation-reversing arclength
isometry. The result is a hyperbolic four-holed sphere with the
indicated separating geodesic. Cutting and canonically capping gives
two marked spheres. On either one, label the external centers $0,1$
and the center of the node cap $\infty$. By
Proposition~\ref{A2-schwarzian}, its projective-connection
coefficients are $a_0=a_1=a_\infty=1$.

Let $u_0$ and $u_1$ be centered unit-disc coordinates on the two
external caps, and use the local coordinates $z$ and $\zeta=1-z$,
respectively. Write
\[
 u_0(z)=c_0z(1+b_0z+O(z^2)),\qquad
 u_1(1-\zeta)=c_1\zeta(1+b_1\zeta+O(\zeta^2)).
\]
The formula in that proposition, first as written and then with the
labels $0,1$ interchanged, gives $b_0=b_1=1/2$. These coefficients
depend only on the two embedded pointed discs: changing either
centered unit-disc coordinate multiplies it by a unit scalar and
does not change $b_i$.

Suppose that after forgetting the cap at $\infty$, the two retained
discs were admissible on the twice-bordered, once-punctured sphere.
Its hyperbolic metric would have a cusp at $\infty$ and positive
boundary lengths $L'_0,L'_1$ at the retained discs. Set
\[
 a'_0=\frac12\left(1+\frac{(L'_0)^2}{4\pi^2}\right)>\frac12,
 \qquad
 a'_1=\frac12\left(1+\frac{(L'_1)^2}{4\pi^2}\right)>\frac12,
 \qquad a'_\infty=\frac12.
\]
The same Schwarzian formula applies when a boundary is a cusp: a
cusp developing coordinate is a constant multiple of $\log q$,
whose Schwarzian is $1/(2q^2)$. Admissibility would identify the
canonical cap domains with the two already given discs. Their
unchanged coefficients would therefore satisfy
\[
 \frac{a'_0+a'_1-1/2}{2a'_0}=\frac12,
 \qquad
 \frac{a'_0+a'_1-1/2}{2a'_1}=\frac12.
\]
The first equality gives $a'_1=1/2$ and the second $a'_0=1/2$,
contradicting both positive boundary lengths. This is an obstruction
to the full disc shapes; it allows the new lengths to vary freely.
\end{proof}

The word ``admissible'' could instead be assigned a new meaning on
a nodal curve, namely joint extendibility by suitable node caps.
That is a possible repair, but it is a different condition from
admissibility of just the external system on the punctured
normalization. The next theorem describes the repaired data and
their image without retaining the node discs as independent
parameters.

Let $L$ be the set of actual external boundary components; other
external markings, if present, remain cusps. For each vertex $v$,
let $J_v$ contain all node branches and all legs in $L$ incident to
$v$. Put $P=E\sqcup L$ and define
\[
 B_v:\R^P\longrightarrow\R^{J_v},\qquad
 (B_vt)_a=
 \begin{cases}
 t_e,&a\text{ is a branch of }e,\\
 t_l,&a\text{ is the external leg }l.
 \end{cases}
\]
In particular a loop edge contributes two entries, while a leg
contributes one. For the fixed marked normalization $Y=(Y_v)$,
let $S_v,K_v$ be the cap potential and support function with active
set $J_v$ and all remaining markings as cusps, and define
\[
 \Psi_{\Gamma,L}(t)=\sum_v S_v(B_vt),\qquad
 K_{\Gamma,L}(p)=\sum_v K_v(B_vp).
\]
An empty active set contributes zero. Choose tangent frames only
to write radii as positive numbers; the conditions below are
intrinsic under changes of those frames.

\begin{theorem}[Proved: exact bordered model by tensors and shadows]
\label{A1-bordered-shadow-range}
On the fixed graph and normalization, record the nonzero node
tensors $q_e$ and the external cap shadows, of radii $r_l>0$.
Canonical cutting and capping gives a real-analytic diffeomorphism
from bordered cut surfaces, with actual toric seam phases, onto
\begin{align*}
 \mathcal V_{\Gamma,L}(Y)=\biggl\{(q,r):\quad&
 \sum_{e\in E}p_e\log|q_e|+
 \sum_{l\in L}p_l\log r_l<K_{\Gamma,L}(p)\\
 &\text{for every }0\ne p\in[0,\infty)^P\biggr\}.
\end{align*}
The logarithmic radial image is open, convex, and downward closed.
The source length at every edge or external leg is
$2\pi\sqrt{t_j}$, where $t$ is the unique solution
\[
 \nabla\Psi_{\Gamma,L}(t)
   =\bigl((\log|q_e|)_{e\in E},(\log r_l)_{l\in L}\bigr).
\]
Thus the node tensors and the external shadows determine all node
and external cap domains, all boundary lengths, and the cut surface.
The statement includes self-nodes and fixed cusp markings. Its
tangent map is an isomorphism; as the normalization varies, the
range forms an open real-analytic subbundle. The inclusion into
the full nonzero-tensor and positive-shadow bundle is an
equivariant homotopy equivalence.
\end{theorem}
\begin{proof}
The chain rule and the seam tensor formula give
\[
 \frac{\partial\Psi_{\Gamma,L}}{\partial t_e}
       =\log r_{a_e}+\log r_{b_e}=\log|q_e|,
 \qquad
 \frac{\partial\Psi_{\Gamma,L}}{\partial t_l}=\log r_l.
\]
By Theorem~\ref{A2b-energy} its Hessian is
\[
 d^2\Psi_{\Gamma,L}=
       \sum_v B_v^T(d^2S_v)B_v>0.
\]
Indeed every coordinate of $\R^P$ occurs in at least one $B_v$,
so their combined map is injective. Integrating this positive
quadratic form along a segment proves global injectivity, and
the inverse function theorem proves local analytic invertibility.

Lemma~\ref{A2b-recession} gives recession function
$K_{\Gamma,L}$. The strict inequalities are necessary by the
componentwise inequalities of Theorem~\ref{A2b-image}. Conversely
let $x$ satisfy the displayed support inequalities. The convex
function $\Psi_{\Gamma,L}-x\cdot t$ has bounded sublevels: an
unbounded sublevel sequence, normalized in direction, would
produce $p\geq0$, $p\ne0$, with recession slope
$K_{\Gamma,L}(p)-p\cdot x\leq0$. Its lower semicontinuous closure
has a minimizer in the closed orthant. No finite minimizer lies
on a boundary face, because the inward diagonal derivative is a
sum of cap log radii with positive multiplicities; a vanishing
coordinate supplies at least one term tending uniformly to
$-\infty$ by the collar estimate, and all remaining terms have
uniform upper bounds. The unique minimizer is therefore interior
and satisfies $\nabla\Psi_{\Gamma,L}=x$. This proves the exact
radial image and its analytic inverse.

For those recovered parameters the canonical cap systems exist
uniquely by Theorem~\ref{A4-grafting}. Edge branches have equal
lengths by construction. The phase of $q_e$ determines their toric
identification, and Proposition~\ref{A1-retained-discs} reconstructs
the surface. The same actual-phase convention as in
Theorem~\ref{A1-graph-range} applies; fully marked real twists
require the corresponding covering on the target.

With varying normalization, the derivative is block triangular
with the identity on normalization parameters, the positive
Hessian on lengths, and the identity on phases. Hence it is an
isomorphism and the inverse varies real-analytically. Finally set
every $t_j=1$. The resulting canonical shadows define a positive
reference radius at each leg and a reference circle in every
tensor line. Linear interpolation of logarithmic radii to these
references is a deformation retraction of both the full radial
bundle and the stated convex subbundle onto the same phase-torus
bundle. It is intrinsic and respects automorphisms, as in
Theorem~\ref{A1-graph-range}. This proves the last assertions.
\end{proof}

\begin{corollary}[Proved: the image when full external discs are retained]
\label{A1-bordered-shape-image}
Replace the manuscript's external-only admissibility requirement
by arbitrary embedded pointed external discs with disjoint
analytic closures avoiding every other special point. The exact
image of canonical surgery in this target consists of the tuples
$(Y,(q_e),(D_l))$ satisfying both of the following conditions:
\begin{enumerate}
\item The tensors and the shadows $r_l=r(D_l)$ satisfy all support
inequalities of Theorem~\ref{A1-bordered-shadow-range}.
\item Let $t$ be the unique gradient inverse furnished by that
theorem. On every component, the given external discs $D_l$ are
exactly the external members of the canonical full cap system
with squared length parameters $B_vt$.
\end{enumerate}
These conditions determine a real-analytic graph over
$\mathcal V_{\Gamma,L}(Y)$. They give a complete reconstruction
criterion, and do not require the retained external subsystem
to be admissible after the node caps have been discarded.
\end{corollary}
\begin{proof}
Both conditions hold for a surgery output. Conversely condition
(1) uniquely constructs the full canonical cap systems, equal
paired lengths, and phase identifications. They weld by
Proposition~\ref{A1-retained-discs}; condition (2) says exactly
that their external members are the prescribed discs. The
constructed surface is therefore a preimage, and uniqueness
follows from Theorem~\ref{A1-bordered-shadow-range}. Analytic
dependence of the reconstructed disc embeddings is
Lemma~\ref{A2b-regularity}. Proposition~\ref{A1-bordered-landing}
shows why external-only admissibility cannot be imposed in
addition as an automatic property.
\end{proof}

\subsection{Toric cuts with unequal paired lengths}
\label{sec:A1-toric-unequal}

There is a fourth reading of the cut datum which must be distinguished
from the geodesic reading. Require that the identification of the two
banks of every cut be toric for their canonical, normalized arclength
circle actions, but do not require their hyperbolic circumferences to
be equal. Toric compatibility is a condition on normalized angles;
it does not equate these two lengths. In this subsection a cut is
an actual analytic embedded seam, and isomorphisms must preserve the
seams and the external labels. The phase is an actual circle
identification, rather than a fully marked real twist.

Fix a stable graph $\Gamma$, its marked normalization $Y=(Y_v)$,
and a subset $L$ of the external markings which are actual boundary
components; the other external markings remain cusps. Let $H$ be
the set of the $2|E|$ node branches together with $L$, and let
$P=E\sqcup L$. Every branch or leg in $H$ has its own independent
positive squared length parameter
\[
 s_h=(\ell_h/2\pi)^2>0.
\]
At each component choose its unique canonical cap system with these
parameters, as supplied by Theorem~\ref{A4-grafting}. Glue the node
boundaries torically. The analytic collar gluing, equivalently the
analytic case of Theorem~\ref{A3-welding}, produces a Riemann surface
with its actual seams even when the paired lengths differ. Conversely,
cutting such a surface and canonically capping recovers these data.
We denote this space, or the corresponding groupoid when
automorphisms are retained, by $\mathcal R_{\Gamma,L}(Y)$.

Put $S(s)=\sum_v S_v(s|_{J_v})$, where $J_v$ is the active set at
the vertex. Theorem~\ref{A2b-image} identifies the full cap logarithm
map
\[
 F=\nabla S:(0,\infty)^H\longrightarrow D_Y
\]
with a real-analytic diffeomorphism onto an open convex downward
closed subset $D_Y\subset\R^H$. A vertex with no active flags
contributes no variables. In fixed tangent frames define the
surjective linear map
\[
 P_*:\R^H\longrightarrow\R^P,\qquad
 (P_*x)_e=x_{h_e^+}+x_{h_e^-},\quad
 (P_*x)_l=x_l.
\]
Thus the logarithms of the node tensor norms and external shadows
are exactly $P_*F(s)$. Changes of frames translate these formulas
compatibly and do not change their intrinsic meaning.

\begin{theorem}[Proved: the broader toric domain, its exact image, and its extra fibers]
\label{A1-toric-domain}
The tensor and external-shadow map from $\mathcal R_{\Gamma,L}(Y)$
has exactly the same image $\mathcal V_{\Gamma,L}(Y)$ as the
equal-length geodesic map of
Theorem~\ref{A1-bordered-shadow-range}; for $L=\varnothing$ this is
$\mathcal U_\Gamma(Y)$ of Theorem~\ref{A1-graph-range}. In particular
\[
 P_*(D_Y)=\left\{y:\ p\cdot y<K_{\Gamma,L}(p)
                   \text{ for all }0\ne p\in[0,\infty)^P\right\}.
\]
This is a proper subset of the unrestricted tensor/shadow fiber
whenever $P\ne\varnothing$; when $P=\varnothing$ it is the
entire one-point fiber.
Before finite automorphism quotients, the fiber over a fixed tensor
and shadow tuple is an open convex set of dimension $|E|$ in its
full cap logarithms. More strongly, after orienting each edge, the
quantities
\[
 d_e=s_{h_e^+}-s_{h_e^-}\in\R
\]
give a global real-analytic identification
\[
 \mathcal R_{\Gamma,L}(Y)
       \simeq\mathcal V_{\Gamma,L}(Y)\times\R^E,
\]
under which the surgery map is projection and the geodesic submodel
is $d=0$. Thus its tangent map is surjective with kernel dimension
$|E|$; it is an isomorphism precisely when $E=\varnothing$.
If the normalization varies and the graph has total genus $g$ and
$n$ external markings, the broader toric source has real dimension
$6g-6+2n+|L|+|E|$, whereas the tensor/shadow target has dimension
$6g-6+2n+|L|$. Thus the dimension count in 67/077 applies to the
geodesic model, not to actual toric cuts of independent lengths.

The identification is intrinsic if $\R^E$ is regarded as the real
space of antisymmetric functions on the two branches of each edge.
Branch interchange changes the corresponding $d_e$ to $-d_e$,
and graph automorphisms also permute edges. The construction varies
real-analytically with the normalization and is equivariant.
Consequently it defines a real vector orbibundle over the geodesic
image when finite automorphisms are retained. Scaling $d$ to zero
is a fiber-preserving equivariant strong deformation retraction.
The broader toric surgery map therefore induces the same
fundamental-group isomorphism to the unrestricted tensor/shadow
target, although its tangent map has a nonzero kernel whenever
there is a cut.
\end{theorem}
\begin{proof}
For a point $x\in D_Y$ and $p\geq0$, the vector $P_*^Tp$ is
nonnegative, and its entries at the two branches of an edge both
equal $p_e$. The component support inequalities give
\[
 p\cdot P_*x=(P_*^Tp)\cdot x<K_{\Gamma,L}(p)
 \qquad(p\ne0).
\]
Conversely, the already proved equal-length gradient inverse in
Theorem~\ref{A1-bordered-shadow-range} lifts every point satisfying
these inequalities to $F(P_*^Tt)\in D_Y$. This proves equality of
the two images without an additional separation or closure
assumption. All tensor phases occur independently by
Lemma~\ref{A1-tensor}. If $P$ is nonempty, a coordinate unit
vector $p$ gives a finite upper bound on the corresponding
logarithmic coordinate, proving properness. The empty case is
the identity on its one-point fiber.

Fix the target tuple, hence its logarithmic vector $y$ and its
tensor phases. Its cap-logarithm fiber is
\[
 D_Y\cap P_*^{-1}(y).
\]
It is nonempty, convex, and open in the affine space
$P_*^{-1}(y)$, whose dimension is
$|H|-|P|=|E|$. The tensor and the two shadow circles uniquely
determine the toric boundary identification, so there are no
additional phase variables in this fiber.

We prove the stronger global parametrization rather than infer it
merely from contractibility. Prescribe any $d\in\R^E$, put
$a(d)_{h_e^\pm}=\pm d_e/2$ and $a(d)_l=0$, and consider
\[
 \Phi_d(t)=S(P_*^Tt+a(d)),\qquad
 \Omega_d=\{t_e>|d_e|/2,\ t_l>0\}.
\]
Its Hessian is $P_*\,d^2S\,P_*^T>0$, because $P_*^T$ is
injective. Its recession function in every nonnegative direction
is $K_{\Gamma,L}$: apply Lemma~\ref{A2b-recession} to the positive
starting vector $P_*^Tt_0+a(d)$ and direction $P_*^Tp$.

For every $y$ in the displayed strict support region,
$\Phi_d(t)-y\cdot t$ has bounded sublevels, by the same
escaping-sequence segment argument as in
Theorem~\ref{A2b-image}. The finite lower bounds defining
$\Omega_d$ ensure that every limiting recession direction is
nonnegative. Its lower semicontinuous closure attains a minimum.
A finite boundary point has at least one branch or external
squared length tending to zero. Along an inward diagonal the
derivative is the sum of the cap logarithms, each with positive
multiplicity; at least one tends uniformly to $-\infty$ by the
collar bound, while every remaining term has a finite upper bound.
Thus no boundary minimum is possible. The minimum is interior
and unique, and satisfies
\[
             \nabla\Phi_d(t)=P_*F(P_*^Tt+a(d))=y.
\]
This constructs exactly one cap system for every $(y,d)$. The
implicit function theorem applied to its positive Hessian gives
joint real-analytic dependence on $y,d$ and the normalization.
There is no nonsmoothness at $d_e=0$: near any solution the
arguments of $S$ remain strictly positive, and its defining
equation is analytic in $(t,d)$; the expression $|d_e|/2$
describes only the open domain of existence.

The unique seam phases prescribed by the tensors complete the
claimed diffeomorphism. If every $d_e=0$, equal circumferences
make the toric seams isometries and their hyperbolic collars glue
smoothly, as in Proposition~\ref{A1-retained-discs}. Conversely,
if every seam is a geodesic for the canonical metric of the welded
surface, every cut component inherits its unique canonical
hyperbolic metric, so the two circumferences agree at every edge.
Thus $d=0$ is exactly the geodesic submodel. This converse uses
all seams simultaneously; a single geodesic among other
nongeodesic cuts does not justify the restriction argument.

Lengths are intrinsic, so any permitted normalization isomorphism
acts on $d$ by the stated signed edge permutation. Uniqueness of
the inverse makes the whole construction equivariant, including
self-loops. The homotopy $(q,r,d)\mapsto(q,r,(1-u)d)$ consequently
descends to the quotient groupoids and their coarse spaces. Its
composition with the common radial retract of
Theorem~\ref{A1-bordered-shadow-range} proves the
fundamental-group assertion. The vector-orbibundle description
does not identify a coarse quotient fiber with a smooth vector
space when a nontrivial finite stabilizer acts on it.
Finally, the marked normalization stratum has real dimension
$6g-6+2n-2|E|$. Adding $2|E|+|L|$ independent cap lengths and
$|E|$ seam phases gives the asserted source dimension; the
geodesic section removes precisely $|E|$ length differences.
\end{proof}

\begin{proposition}[Refuted: toric compatibility alone implies an inverse]
\label{A1-toric-noninjective}
Even for the unbordered four-labelled sphere with one separating
analytic cut, the torically compatible surgery map is not
injective. There are distinct such cut objects with fixed
normalization and tensor $q=\partial_z|_0\otimes\partial_w|_0$.
\end{proposition}
\begin{proof}
Let $P(r)$, $0<r<4$, be the uniquely determined marked hyperbolic
pant with two cusps and one geodesic boundary whose canonical cap,
on its normalized sphere $(0;1,\infty)$, has conformal radius $r$
at $0$. Existence, uniqueness, and strict variation of its
boundary length $\ell(r)$ follow from the one-active-cap case of
Theorem~\ref{A2b-image} and the Koebe endpoint in
Corollary~\ref{A2-explicit-limits}. Normalize its cap embedding
$f_r:\overline\D\to\mathbf P^1$ by $f_r(0)=0$ and
$f_r'(0)=r>0$. This removes the rotational ambiguity.

For each of the following two ordered pairs of pants, glue their
boundaries by the explicit toric rule
\[
 f_r(e^{i\theta})\sim f_{r'}(e^{-i\theta}):
 \qquad (r,r')=(1,1),\qquad (r,r')=(2,1/2).
\]
The cap maps extend holomorphically across their analytic boundaries,
so the rule gives analytic conformal welding. Label the two cusp
points on each side separately, so that any isomorphism of cut
objects preserves the two sides. Both outputs have normalization
equal to two spheres with their node at $0$ and the labelled
markings at $1,\infty$. By Lemma~\ref{A1-tensor}, their tensor is
in both cases $rr'\,\partial_z\otimes\partial_w
=\partial_z\otimes\partial_w$.

They are nevertheless distinct cut objects: a conformal
isomorphism preserving the cut restricts to a conformal isomorphism
of the first cut components, and therefore preserves their unique
canonical hyperbolic boundary lengths. Those lengths are
$\ell(1)$ and $\ell(2)$, which are unequal by strict monotonicity.
The first seam is geodesic. The second cannot be geodesic, since
its two canonical cut-component lengths $\ell(2)$ and $\ell(1/2)$
are unequal. This argument distinguishes the actual cut objects;
it does not assume that forgetting their cuts gives distinct
underlying four-pointed surfaces.
\end{proof}

\begin{corollary}[Proved: what changes when full external discs are retained]
\label{A1-toric-full-discs}
For the broader toric interpretation, retain arbitrary pointed
external discs with disjoint analytic closures, in addition to
the tensors. Such a tuple $(Y,q,(D_l))$ is a surgery output if and
only if its external shadows $r_l=r(D_l)$ and tensors belong to
$\mathcal V_{\Gamma,L}(Y)$ and there exists $d\in\R^E$ for which
the unique canonical cap system constructed in
Theorem~\ref{A1-toric-domain} has external members exactly $D_l$.
Its fiber is precisely this set of $d$'s, with the prescribed
finite automorphism identifications. In particular the
$\R^E$-fiber assertion concerns the tensor/shadow target; it
must not be transferred to a target retaining full external
disc shapes. The geodesic criterion of
Corollary~\ref{A1-bordered-shape-image} is the special case $d=0$.
\end{corollary}
\begin{proof}
For fixed $(Y,q,r,d)$ the preceding theorem constructs all caps and
their toric identifications uniquely. Necessity follows by cutting
a source object and recording its individual squared lengths.
Conversely the stated equality of external members makes the
constructed source map to exactly the prescribed full-disc tuple.
No claim that the external subsystem remains admissible after
discarding the node caps is used; that distinct condition still
has the obstruction in Proposition~\ref{A1-bordered-landing}.
\end{proof}

\subsection{Raw toric surgery does not descend to cut isotopy classes}
\label{sec:A1-toric-descent}

There remains a distinction between replacing a cut by its unique
geodesic and asserting that raw cutting and capping is independent
of its representative. The latter is suggested by the discussion
of a fixed welded $X$ and nearby seams in 67/079--080. It is false
even if all the representatives under consideration are torically
compatible. We give an exact analytic counterexample; its phase
is specified by a deterministic analytic rule, not by an evaluated
numerical angle.

\begin{lemma}[Proved: the third jet of a one-active-cap pant]
\label{A1-pant-third-jet}
Let the hyperbolic pant have one boundary of length $2\pi\lambda$,
$\lambda>0$, and two cusps. Normalize its capped sphere by placing
the cap centre at $0$ and the ordered cusps at $1,\infty$.
Let $r(\lambda)$ be the cap's conformal radius, and normalize its
unit-disc coordinate $u_\lambda$ to have positive derivative at
zero. Then
\[
 g_\lambda(z):=r(\lambda)u_\lambda(z)
       =z+\frac12z^2+\beta(\lambda)z^3+O(z^4),
 \qquad
 \beta(\lambda)=\frac5{16}
                 +\frac1{16(\lambda^2+4)}.
\]
Moreover $r:(0,\infty)\to(0,4)$ is a strictly increasing
real-analytic bijection.
\end{lemma}
\begin{proof}
The radius assertion is the one-active-cap case of
Theorem~\ref{A2b-image} and the Koebe endpoint in
Corollary~\ref{A2-explicit-limits}. Put
$a=(1+\lambda^2)/2$. We use the projective-connection argument
of Proposition~\ref{A2-schwarzian} with its cusp extension made
explicit. At a cusp a developing coordinate is a nonzero constant
multiple of $\log q$, whose Schwarzian is $1/(2q^2)$.
Thus the double-pole coefficients at $1,\infty$ are both $1/2$.
At the positive boundary they are $(1+\lambda^2)/2$ as before.
The unique projective connection on the three-pointed sphere with
these principal coefficients is consequently
\[
 T(z)=\frac a{z^2}+\frac1{2(z-1)^2}
                         -\frac a{z(z-1)}
      =\frac a{z^2}+\frac a z+(a+1/2)+O(z).
\]
The cap developing map gives
$T=a(g'/g)^2+\{g,z\}$. For
$g=z+bz^2+\beta z^3+O(z^4)$, its coefficients of $z^{-1}$ and
$z^0$ are respectively
\[
 2ab,\qquad (4a+6)\beta-(a+6)b^2.
\]
Matching the displayed connection gives $b=1/2$ and
\[
 \beta=\frac{5a+8}{16a+24}
       =\frac5{16}+\frac1{16(\lambda^2+4)}.
\]
The normalization by $r$ changes neither logarithmic derivatives
nor the Schwarzian, so this calculation applies to the stated
normalized germ.
\end{proof}

\begin{lemma}[Proved: a four-point plumbing coefficient]
\label{A1-four-point-plumbing-jet}
On two spheres with node $0$ and external markings $1,\infty$,
take node-coordinate germs
\[
 g_i(z)=z+\frac12z^2+\beta_i z^3+O(z^4),\qquad i=1,2.
\]
Plumb by $g_1(z)g_2(w)=t$. On the resulting four-pointed sphere,
normalize the global coordinate so that the first component's
$\infty$ is $\infty$, its $1$ is $1$, and the second component's
$\infty$ is $0$. Let $\Lambda(t)$ be the coordinate of the second
component's $1$. Then
\[
 \Lambda(t)=t-\frac12t^2+
 \left(\frac9{16}-\frac54(\beta_1+\beta_2)
                         +4\beta_1\beta_2\right)t^3+O(t^4).
\]
In particular, for the germs of Lemma~\ref{A1-pant-third-jet},
the cubic coefficient is
\[
 C(\lambda_1,\lambda_2)
   =\frac{11}{64}
       +\frac1{64(\lambda_1^2+4)(\lambda_2^2+4)}.
\]
\end{lemma}
\begin{proof}
For this coefficient calculation choose sufficiently small fixed
plumbing discs whose closures lie in $|z|<1/2$ and $|w|<1/2$.
This is possible for any of the coordinate germs in the statement.
The retained pieces then contain the circles $|z|=1$, $|w|=1$
and both external markings. First normalize a meromorphic
coordinate $Z$ to have its sole simple
pole at the first $\infty$, leading term $z$ there, and value zero
at the second $\infty$. On the two retained pieces write
\[
 Z_1(z,t)=z+B(t)+\sum_{n\geq1}C_n(t)z^{-n},\qquad
 Z_2(w,t)=\sum_{n\geq1}D_n(t)w^{-n}.
\]
These are convergent Laurent expansions for sufficiently small
$t$, in particular at $z=w=1$; their coefficients are meromorphic
on the whole central components. They depend holomorphically
on $t$, including $t=0$, where
$Z_1=z$ and $Z_2=0$. One can see this directly from the family of
meromorphic functions with the specified simple pole: on the
central genus-zero nodal curve its pole divisor has multidegree
$(1,0)$, with vanishing first cohomology, and evaluation and the
leading coefficient give the two normalizations. Equivalently,
Laurent principal-part matching in the plumbing annulus constructs
the same functions.

In the local total space $xy=t$, holomorphicity at the node and
$Z_2(w,0)=0$ show that
\[
 C_n(t)=O(t^{n+1}),\qquad D_n(t)=O(t^n).
\]
Indeed a term $y^m$ on the second branch has coefficient vanishing
at least once in $t$; after $y=t/x$ it contributes at least
$t^{m+1}$ to a pole of order at most $m$ on the first branch.
The analogous first-branch terms give the stated bound for $D_n$.

For completeness all terms needed for the cubic coefficient are
as follows. Set $A_i=\beta_i-1/4$ and $a=1/2-\beta_1$.
The inverse germ gives, on the overlap,
\[
 z=\frac t{g_2(w)}-\frac12\frac{t^2}{g_2(w)^2}
                          +a\frac{t^3}{g_2(w)^3}+O(t^4).
\]
Matching negative and constant Laurent coefficients at $w=0$
gives
\begin{align*}
 C_1(t)&=A_2t^2+O(t^3),\\
 B(t)&=\tfrac12t+(\tfrac12-\tfrac32\beta_2)t^2+O(t^3),\\
 D_1(t)&=t+\tfrac12t^2+
       \left[a(\tfrac32-3\beta_2)+A_1A_2\right]t^3+O(t^4),\\
 D_2(t)&=-\tfrac12t^2-\tfrac32a\,t^3+O(t^4),\\
 D_3(t)&=a\,t^3+O(t^4).
\end{align*}
For example, the coefficient of $w$ in $1/g_2(w)$ is
$1/4-\beta_2=-A_2$, which gives the first line. The negative
Laurent parts of $g_2^{-2}$ and $g_2^{-3}$ give the next three
principal-part lines; the $A_1A_2$ term comes from the coefficient
of $t/g_2(w)$ in $C_1(t)/z$. Higher $C_n$ cannot contribute to
these negative parts before order $t^4$.

It follows that
\begin{align*}
 Z_1(1,t)&=1+\tfrac12t+(\tfrac14-\tfrac12\beta_2)t^2+O(t^3),\\
 Z_2(1,t)&=t+\left[a(1-3\beta_2)+A_1A_2\right]t^3+O(t^4).
\end{align*}
The required cross-ratio is $\Lambda=Z_2(1,t)/Z_1(1,t)$.
Division gives the stated symmetric cubic coefficient. Finally
\[
 \frac9{16}-\frac54(\beta_1+\beta_2)+4\beta_1\beta_2
 =\frac{11}{64}
      +4(\beta_1-5/16)(\beta_2-5/16),
\]
which gives the last formula. This is a coefficient identity
derived from convergent plumbing, not a numerical approximation.
Enlarging the coordinate discs gives the same family wherever
both constructions are defined, by
Proposition~\ref{A1-welding-plumbing}; thus later larger parameter
domains give its holomorphic continuation and do not change
these Taylor coefficients.
\end{proof}

\begin{theorem}[Refuted: raw toric surgery is invariant under isotopy in a fixed surface]
\label{A1-toric-nondescent}
There is a four-labelled punctured sphere $X$ with two isotopic
essential analytic cuts $C,C_{\mathrm{geo}}$ such that both cuts
are canonically torically compatible, $C_{\mathrm{geo}}$ is the
unique hyperbolic geodesic in their isotopy class, and raw canonical
cutting and capping produces different nodal tensors. Thus raw
toric surgery does not descend to a functor on cut isotopy classes.
Replacing a representative by its geodesic before applying surgery
defines the different, well-defined functor already studied.
\end{theorem}
\begin{proof}
We specify all the choices. Use the radius function in
Lemma~\ref{A1-pant-third-jet}, and set
\[
 r_1=r(1),\qquad r_*=r(1/2),\qquad
 r_2=\frac{r_*^2}{4r_1},\qquad
 r_0=\frac{r_*}{2},\qquad R=r_0^2=r_1r_2.
\]
These are positive and less than $4$. Let
$\lambda_2=r^{-1}(r_2)$ and $\lambda_0=r^{-1}(r_0)$; strict
monotonicity gives $0<\lambda_2<\lambda_0<1/2$.
The reference pants defining $r_1,r_*$ have the explicit boundary
lengths $2\pi,\pi$. All other pants are uniquely fixed by the
displayed conformal radii, so no unchosen cap realization remains.

For the fixed pair $(\lambda_1,\lambda_2)=(1,\lambda_2)$ let
$\Lambda_U(t)$ be the cross-ratio of its plumbing family from the
preceding lemma. For the fixed equal pair $(\lambda_0,\lambda_0)$
write $\Lambda_G(t)$. These are honest holomorphic families on a
disc containing $|t|=R$, not just formal germs. To make this
domain explicit, put
\[
 \delta(\lambda)=\arctan\!\bigl(\operatorname{csch}(\pi\lambda)\bigr),
 \qquad
 a(\lambda)=r(\lambda)
                 \exp\!\left(\frac{\delta(\lambda)}{2\lambda}\right).
\]
The collar extension used in Theorem~\ref{A2b-image} makes
$g_\lambda$ a valid node coordinate on an embedded disc of radius
$a(\lambda)>r(\lambda)$, avoiding both external markings.
The factor $1/2$ keeps this disc strictly inside the available
collar. Thus the usual plumbing quotients exist for
\[
 |t|<\min\{a(1)a(\lambda_2),a(\lambda_0)^2\}=R\rho,
 \qquad \rho>1.
\]
At $t=R\eta$, $|\eta|=1$, their cut along the circles
$|g_i|=r_i$ is precisely the actual toric welding
$u_1u_2=\eta$, by Proposition~\ref{A1-welding-plumbing}.

Define the holomorphic function on $|c|<\rho$
\[
                 \Delta(c)=\Lambda_U(Rc)-\Lambda_G(Rc).
\]
Its cubic coefficient is nonzero. Indeed, with
$h(\lambda)=1/(\lambda^2+4)$,
\[
 h(1)h(\lambda_2)<\frac1{20}
       <\frac{16}{289}<h(\lambda_0)^2,
\]
and Lemma~\ref{A1-four-point-plumbing-jet} gives
\[
 [c^3]\Delta(c)
   =\frac{R^3}{64}
       \left[h(1)h(\lambda_2)-h(\lambda_0)^2\right]
   <-\frac{31R^3}{369920}<0.
\]
Consequently $\Delta$ is not the zero function. It has only
finitely many zeros on the unit circle, since that circle is a
compact subset of its open holomorphy disc.

For a deterministic phase, define
\[
 N=\min\{n\geq1:\ \Delta(e^{i/n})\ne0\},
 \qquad \eta_*=e^{i/N}.
\]
This integer exists: otherwise the distinct zeros $e^{i/n}$ would
accumulate at $1$, an interior point of the holomorphy disc, and
the identity theorem would contradict the nonzero cubic coefficient.
This is an exact analytic specification of the phase; no numerical
value of $N$ is asserted or needed. More generally every unit
phase outside the finite zero set gives the same counterexample
construction.

Now take $X$ to be the unequal-pair welding
$P(r_1)\cup_{u_1u_2=\eta_*}P(r_2)$, with its actual analytic
seam $C$ and its four distinct external labels. Its normalized
nodal output is the fixed pair of three-pointed spheres, with
tensor $q=R\eta_*$. The equal-pair welding
$X_G=P(r_0)\cup_{u_1u_2=\eta_*}P(r_0)$ has the same tensor,
and its seam is geodesic. But
$\Delta(\eta_*)\ne0$ says that $X$ and $X_G$ are not isomorphic
as four-labelled punctured spheres: their specified cross-ratios
differ.

Let $C_{\mathrm{geo}}$ be the unique hyperbolic geodesic of $X$
in the essential isotopy class of $C$. Raw cutting along this
geodesic is torically compatible. If its tensor were also
$R\eta_*$, the unique geodesic inverse in
Theorem~\ref{A1-graph-range}, or its explicit four-point case
Corollary~\ref{A1-small-ranges}, would identify its welded surface
with $X_G$. This contradicts the cross-ratio inequality. Thus
the two isotopic cuts in this same $X$ have different tensors,
which proves non-descent. No inference from the dimension of
the fixed-target fibers is used.
\end{proof}

\subsection{The algebraic-category obstruction in the four-pointed case}

\begin{proposition}[Refuted: finite-type algebraicity of the natural cut component]
\label{A1-cut-not-algebraic}
Let $\mathcal D$ be a connected component of the moduli space of
four-labelled punctured spheres equipped with one essential
unoriented cut isotopy class, represented by its hyperbolic
geodesic. Give $\mathcal D$ its natural complex structure: forgetting
the locally constant cut class is a local biholomorphism to
$M_{0,4}^{\mathrm{an}}$.
Then
\[
 \mathcal D\simeq
 \mathbb H/\langle\tau\mapsto\tau+2\rangle
 \simeq\{q\in\C:0<|q|<1\}.
\]
In particular $\mathcal D$ is not the analytification of a smooth
connected finite-type algebraic curve over $\C$. Thus the natural
geodesic-cut source cannot be treated as a finite-type algebraic
moduli curve, and it cannot be algebraically isomorphic to the
corresponding unrestricted nodal-tensor stratum $\C^\times$.
This is a limitation of that specified algebraic interpretation;
it does not infer nonholomorphicity of the surgery map from its
real-analytic description.
\end{proposition}
\begin{proof}
We give the passage from the elliptic double cover explicitly.
The double cover of a four-pointed sphere branched at its four
markings has genus one by Riemann--Hurwitz. Choose one ramification
point as origin. The other three are its nonzero two-torsion
points. Conversely an elliptic curve divided by $z\mapsto-z$
recovers the sphere and these four markings. Markings lift in two
ways, differing by $-1$. Thus the effective marked parameter is
the period ratio $\tau\in\mathbb H$ for
$E_\tau=\C/(\Z+\tau\Z)$. This is a holomorphic identification
of the Teichm\"uller spaces: locally the double-cover family has
equation $y^2=x(x-1)(x-\lambda)$, and the ratio of the two marked
periods of $dx/y$ varies holomorphically; the inverse is the
quotient elliptic family. The torus parameter itself is
\cite[Example 2.3, p.~12]{AHubbardKoch2014}.

The lifting and curve correspondence are proved in
\cite[Propositions 2.6--2.7, pp.~55--57]{FarbMargalit2012A567}.
Here is the additional restriction needed for labelled markings.
Fixing the chosen origin removes the two-torsion translations;
fixing all four labels requires the homology matrix to be the
identity modulo two. The effective pure group is consequently
$\Gamma(2)/\{\pm I\}$. Under the double cover an essential sphere
curve lifts to two parallel primitive torus curves. Choose a
homology basis in which their common unoriented class is
$\pm(1,0)$. Its stabilizer in the pure group consists exactly of
the matrices
\[
 \pm\begin{pmatrix}1&2m\\0&1\end{pmatrix},\qquad m\in\Z.
\]
Indeed preserving that primitive line forces $c=0$ and $a=d=\pm1$,
and congruence modulo two forces $b$ even. A Dehn twist downstairs
lifts to the twists around the two parallel curves, and is the
generator with $|m|=1$. Therefore forgetting the full marking but
retaining the cut gives precisely the claimed cyclic quotient.
The exponential $q=\exp(\pi i\tau)$ identifies it biholomorphically
with the punctured unit disc. This also agrees with the cyclic
parabolic quotient in \cite[Example 7.2, pp.~25--26]{AHubbardKoch2014},
with period two instead of period one. Since an automorphism of a
sphere fixing its four labelled points is the identity, no residual
orbifold stabilizer changes this calculation.

Suppose now that this punctured disc were $X^{\mathrm{an}}$ for a
smooth connected finite-type algebraic curve $X/\C$. Its smooth
projective completion has some genus $g$, and its complement is a
finite nonempty set of $r$ points; nonemptiness follows because
the punctured disc is not compact. The fundamental group of the
resulting punctured surface is free of rank $2g+r-1$. Since the
punctured disc has fundamental group $\Z$, we have $2g+r=2$,
and $r\ge1$ forces $g=0,r=2$. Thus $X\simeq\mathbf P^1\setminus
\{0,\infty\}=\C^\times$. But $\C^\times$ has universal cover
$\C$, whereas the punctured disc has universal cover $\mathbb H$;
these simply connected Riemann surfaces are not biholomorphic.
For example, after identifying $\mathbb H$ with $\D$, a
biholomorphism $\C\to\mathbb H$ would contradict Liouville's
theorem. This proves the contradiction.

The fixed labelled one-node normalization in the target has no
moduli or automorphisms, and its nonzero tangent tensor is an
arbitrary element of a one-dimensional complex vector space.
Hence its unrestricted target stratum is exactly $\C^\times$.
The stated algebraic isomorphism is therefore impossible for the
natural source just analyzed, independently of the cap-size
obstruction and without any assertion about whether a particular
real-analytic formula is holomorphic.
\end{proof}

Both the punctured disc and $\C^\times$ have fundamental group
$\Z$. The local real tangent isomorphism and the fundamental-group
isomorphism in Theorem~\ref{A1-graph-range} are thus fully
compatible with the failure of a global isomorphism onto the
unrestricted target. They cannot substitute for the range
condition or for the algebraic-category hypothesis.

\subsection{What local plumbing does prove}
Hubbard--Koch, Theorem 9.11 \cite{AHubbardKoch2014}, establishes a homeomorphism onto an open neighborhood for sufficiently small plumbing parameters; the sentence immediately following the theorem explicitly says that the plumbing map is not globally injective. Their Theorem 10.1 gives the complex manifold structure and the universal marked family. Their Section 13 compares the quotient with the analytification of the Deligne--Mumford compactification. These statements do not provide an inverse to the unrestricted-tensor map refuted above.

\begin{proposition}[Known with derivation: the intrinsic first-order parameter]\label{A1-first-order}
For an ordinary node with branches $a,b$ on a smooth normalization $Y$, the one-dimensional space of local first-order smoothings is canonically $T_aY\otimes T_bY$. Given analytic coordinates $x,y$ at the branches, the family $xy=t$ represents this identification by $t\,\partial_x\otimes\partial_y$. Coordinate choices enter at higher order. For a fixed stable complex curve, sufficiently small such parameters together with deformation parameters for the normalized pointed components give local analytic plumbing charts, modulo its finite automorphism group.
\end{proposition}
\begin{proof}
For $A=\C[[x,y]]/(xy)$, the hypersurface presentation gives
\[
\operatorname{Ext}^1_A(\Omega_A,A)=A/(x,y)=\C.
\]
Explicitly, the relation $y\,dx+x\,dy=0$ yields the cokernel of $(u,v)\mapsto yu+xv$. Replacing $x,y$ by coordinates with linear terms $ax,by$ replaces $t$ by $ab\,t$, while the tensor frame changes by $(ab)^{-1}$. Thus the tensor is intrinsic. Nonlinear coordinate terms are irrelevant modulo the square of the deformation ideal, but can change actual finite plumbing parameters. The asserted local charts and finite quotient are exactly the application of Hubbard--Koch, Sections 8--10, to the normalization with its branch markings; stability makes the automorphism group finite. This is a local analytic assertion and supplies no global algebraic isomorphism.
\end{proof}


\subsection{The exact comparison with coordinate plumbing}

\begin{proposition}[Proved: toric welding is plumbing with the retained coordinate germs]\label{A1-welding-plumbing}
For the corrected surgery data of Proposition~\ref{A1-retained-discs}, every seam is an ordinary analytic plumbing, using the coordinate germs supplied by its two pointed caps. This assertion identifies the welded surface itself. It does not identify the finite plumbing parameter, in a fixed independent choice of germs, with an arbitrary element of the whole normal tensor line.
\end{proposition}
\begin{proof}
Choose centred unit-disc coordinates $u,v$ on the two caps, so the toric boundary identification is $uv=c$, $|c|=1$. Analyticity of the boundaries extends both coordinates to collars of $|u|=|v|=1$. Choose $R>1$ close enough to one that these extensions are embedded and avoid all other caps. In this collar the welding identifies the annuli $R^{-1}<|u|<R$ and $R^{-1}<|v|<R$ by $v=c/u$. The identity map on the remaining component surfaces therefore identifies the welded surface with the plumbing quotient. This works equally for two branches on the same normalization component because their extended cap neighborhoods are disjoint.

More explicitly set $x=a u$, $y=b v$, for any nonzero constants $a,b$. The plumbing equation is $xy=t$ with $t=ab c$. If $z,w$ are reference coordinates at the two centres and $f=u^{-1}$, $g=v^{-1}$, its normal tensor is
\[
 t\,\partial_x\otimes\partial_y
 =ab c\,\frac{f'(0)}a\frac{g'(0)}b
            \partial_z\otimes\partial_w
 =c f'(0)g'(0)\,\partial_z\otimes\partial_w.
\]
This is exactly the tensor of the toric surgery. Shrinking or enlarging the overlapping annuli leaves the quotient unchanged: extend $u$ to the opposite chart by $c/v$ and identify the overlaps using that same holomorphic function. These maps agree where both are defined, so they glue to inverse biholomorphisms of the two quotients. This direct comparison is also the argument in Kra, Appendix III, \S\S14.2--14.5 \cite{AKra1990}.

The coordinate germs in this construction are part of the recovered cap system and vary with its lengths. Proposition~\ref{A2-schwarzian} shows already in a pant that their second jets vary when the other lengths change. Thus retaining these germs, choosing fixed horocyclic germs, and retaining only first-order normal tensors are three distinct operations. The displayed tensor identity describes their first-order comparison, not an extension of a fixed plumbing chart to all tensors.
\end{proof}

\begin{proposition}[Known with derivation: the scope of Kra's coordinates]\label{A1-kra-comparison}
Fix a maximal pants graph of type $(g,n)$ and put $d=3g-3+n>0$. Kra's horocyclic plumbing parameters range in a domain $D$ satisfying
\[
 \mathbb D_{e^{-2\pi}}^{d}\subset D
                   \subset\mathbb D_{e^{-\pi/2}}^{d}.
\]
They cover the moduli of surfaces of that type, generally many times. On marked Teichm\"uller space the logarithmic coordinates $\tau_j$ are global holomorphic coordinates on their image; the corresponding Dehn twist sends $\tau_j$ to $\tau_j+2$. After quotienting these twists, $t_j=e^{\pi i\tau_j}$ are the plumbing coordinates on $D\cap(\C^\times)^d$. They are not coordinates on an unrestricted normal-tensor bundle.
\end{proposition}
\begin{proof}
These are Theorems~1(a),(b) and~3(a),(b) on pp.~501--503 of \cite{AKra1990}, and the identification following Theorem~3(d). The exponential converts translation by $2$ into identity and converts the stated imaginary-part bounds into the displayed polydisc bounds. Zero parameters describe nodes. The bounded domain and the residual moduli identifications exclude the proposed unrestricted global inverse. Proposition~\ref{A1-welding-plumbing} specifies the comparison when the same full coordinate germs are used.
\end{proof}

\begin{proposition}[Known with derivation: Wolpert's metric comparison]\label{A1-wolpert-comparison}
The local family $xy=t$, $0<|t|<1$, is the annulus $|t|<|x|<1$. Its complete curvature-minus-one metric and core length are
\[
 ds_t^2=\left[
 \frac{\pi}{|\log|t||}
 \csc\!\left(\frac{\pi\log|x|}{\log|t|}\right)
 \left|\frac{dx}{x}\right|\right]^2,
 \qquad L_t=\frac{2\pi^2}{-\log|t|}.
\]
On compact subsets of either punctured limiting branch the metric tends to the cusp metric. For a degenerating family of compact stable curves, Wolpert's Expansion~4.2 compares the global hyperbolic metric with the interpolated cusp/annulus metric; its constants are local in the normalization parameters and coordinate germs. This supplies a metric comparison for plumbing, not a global inverse on normal tensors.
\end{proposition}
\begin{proof}
The formula is \cite[\S0.1, p.~417]{AWolpert1990}. Writing $x=e^{\xi+i\theta}$ gives the hyperbolic strip metric; reflection about $\xi=\tfrac12\log|t|$ fixes its core, and integrating there for $0\le\theta\le2\pi$ gives $L_t$. The expansion $a/\sin a\to1$ gives the cusp limit. Expansion~4.2 and Example~4.3, pp.~445--447, apply after fixing small deformation parameters and plumbing germs. The normalization factors $F'_j(0),G'_j(0)$ rescale the parameters in that source, exactly as in Proposition~\ref{A1-first-order}. Its prescribed-curvature correction is essential for the global surface metric. The complete annulus formula alone does not assert equality with that global metric. The compact-family hypothesis is retained here.
\end{proof}

\begin{proposition}[Known with an exact counterexample: arbitrary plumbing need not be a chart]\label{A1-hinich-counterexample}
Even for $(g,n)=(0,4)$, disjoint coordinate discs and the equation $xy=t$ do not ensure that the plumbing map is locally invertible at every nonzero parameter. An instance of Hinich's counterexample is obtained by
\[
 R=\tfrac32,\quad B=\tfrac34,\quad C=10,\quad
 A=\tfrac{420}{569},\quad \varepsilon=\tfrac1{10},
 \qquad \phi(z)=(z-R)^2.
\]
The critical plumbing parameter is $t=9/11$.
\end{proposition}
\begin{proof}
This is an exact specialization of \cite[\S\S2.1--2.3]{AHinich2010}; we give a direct cross-ratio verification. Let $P_1=\mathbf P^1_z$ have markings $A,B$ and node branch $\infty$, with coordinate $y=9/(10z)$ at that branch. Let $P_2=\mathbf P^1_w$ have markings $c=(C-R)^2=289/4,\infty$ and node branch $R^2=9/4$. On its node disc use
\[
 x=\frac{10}{11}(R-\sqrt w),\qquad
 w=(R-\tfrac{11}{10}x)^2,\quad |x|<1,
\]
with the branch $\sqrt{R^2}=R$. This disc is embedded: the map $\phi$ is injective on $|z|<11/10$, since equal values require $z=w$ or $z+w=3$. The latter is impossible there. Its image avoids $c,\infty$ because both preimages $C,3-C$ lie outside that disc. The other node disc $|y|<1$ avoids $A,B$, since both have modulus less than $9/10$. Thus all required coordinate discs and markings satisfy the plumbing hypotheses.

For $0<|t|<1$ plumb by $xy=t$. Set $k=11t/9$. The global coordinate on the resulting sphere is $w$ on the second piece and $(R-kz)^2$ on the first: these agree on the annular overlap. The first retained piece is $|z|<(9/10)/|t|$, hence $|kz|<11/10$ and its map is injective. Its image is $\phi(\mathbb D_{11/10})$. The second retained image is the complement of $\phi(\overline{\mathbb D}_{(11/10)|t|})$. Their intersection is exactly $\phi(\{(11/10)|t|<|u|<11/10\})$, the plumbing overlap. Thus their images cover $\mathbf P^1$ and they glue to a biholomorphism. The four markings are
\[
 a(t)=(R-kA)^2,\quad b(t)=(R-kB)^2,\quad c=289/4,\quad\infty.
\]
They stay distinct for $0<|t|<1$. A moduli coordinate is
\[
 \lambda(t)=\frac{a(t)-c}{b(t)-c},\qquad
 \lambda'(9/11)=0,\qquad
 \lambda''(9/11)=\frac{271040}{371353867}\ne0.
\]
Consequently the plumbing map ramifies at $9/11$. All rational inequalities, derivatives and the equivalent residue cancellation are checked in \path{code/part_a_plumbing_checks.py}. Our explicit branch assignment also specifies which side receives each cap; it avoids relying on the reversed inside/outside cap wording in the arXiv v2 paragraph on p.~4. Hinich's Corollary~1.1 retains local coordinates sufficiently near the nodal origin and does not contradict this example.
\end{proof}

For Hubbard--Koch the comparison is the universal marked family of \cite[Theorems 9.11 and 10.1, and \S13]{AHubbardKoch2014}: a sufficiently small coordinate plumbing is its pullback by a local analytic classifying map. At a node the map on the transverse tangent lines is the tensor identification in Proposition~\ref{A1-first-order}; using the same component deformations supplies the remaining tangent directions. Finite automorphism quotients must still be taken. Theorem~9.11 is explicitly local, and its following remark excludes global injectivity. Thus neither that theorem nor the other opened plumbing sources repairs the unrestricted target of L2.

\begin{proposition}[Known with derivation: the opened Earle--Sipe coordinate comparison]\label{A1-Earle-Sipe-comparison}
Let $z,w$ be disjoint local unit-disc coordinates on the two branches to be joined. For $0<r_1,r_2<1$ and $t=r_1r_2e^{i\theta}$, plumbing by $zw=t$ is precisely the analytic boundary welding which removes $|z|<r_1$ and $|w|<r_2$ and identifies
\[
 z=r_1e^{i\varphi},\qquad w=r_2e^{i(\theta-\varphi)}.
\]
Its normal tensor is $t\,\partial_z\otimes\partial_w$. This is the same coordinate-level comparison as the original Earle--Sipe construction, with the full local coordinates retained.
\end{proposition}
\begin{proof}
In \cite[\S2, equations (2.1)--(2.2), pp.~80--82]{EarleSipe1991A567}, the two local identifications are $z_1w_1=t$ and $z_2w_2=t^{2p-3}$. Section~6, equations (6.1)--(6.2), p.~88, constructs the identical fibers by removing discs of radii $|t|^{1/2}$ and $|t|^{(2p-3)/2}$, respectively, and identifying their boundary circles by those same equations. These passages are read in the original published article.

For the displayed general radii, set $u=z/r_1$, $v=w/r_2$. The boundary equation becomes $uv=e^{i\theta}$; extending across the analytic circles gives exactly $w=t/z$ on the annular overlap. Moving the cut within that overlap changes neither quotient, by the identity on the retained charts. Finally
\[
 e^{i\theta}\,\partial_u\otimes\partial_v
 =e^{i\theta}r_1r_2\,\partial_z\otimes\partial_w
 =t\,\partial_z\otimes\partial_w.
\]
Thus this opened original supplies a direct comparison with Proposition~\ref{A1-welding-plumbing}. It does not remove the dependence on the coordinate germs or provide an inverse on unrestricted normal tensors.
\end{proof}

\paragraph{What is and is not checked in the joint Earle--Marden source.}
The distinct Earle--Marden chapter \emph{On holomorphic families of Riemann surfaces}, Contemp. Math.~573 (2012), 67--97, must not be substituted for \emph{Holomorphic plumbing coordinates}, Contemp. Math.~575 (2012), 41--52. The latter is exactly reference [24] in the opened \cite{Matsumoto2019A567}. His printed p.~27 reports its \S13, Theorem~II, distinguishes the asserted embedding over a whole Teichm\"uller stratum from Hubbard--Koch's local product, and states a proof-status limitation beyond the maximal-pants case. That is a precisely located later account. The unavailable joint chapter itself has not been directly verified, and the Earle--Sipe proposition above is not relabelled as a theorem from it.

\begin{proposition}[Known with derivation: identification with the opened joint Earle--Marden original]\label{A1-Earle-Marden-original}
The joint manuscript of Earle and Marden,
\emph{Holomorphic Plumbing Coordinates on Teichm\"uller and Compactified
Moduli Space}, dated \emph{Draft 8/24/11}, gives the following identification.
It is an original author manuscript, recovered from the public archive of
Marden's Minnesota papers directory, rather than the distinct published
2012 chapter. The relevant locators are \S4.1.1--4.1.2, pp.~39--40;
\S4.5, p.~50; and \S4.8.1--4.8.2, equations~(4.8.1)--(4.8.2),
p.~57 \cite{EarleMarden2011DraftA567}.

Normalize the peripheral translation to $A(\zeta)=\zeta+1$.
Suppose the two disjoint puncture neighborhoods are the quotients of
$\operatorname{Im}\zeta>h$ and $\operatorname{Im}\eta<-h$, where $h>0$.
For $\tau=a+2si$, $s>h$, the identification
$\zeta=\eta+\tau$ becomes classical plumbing
\[
 z=e^{2\pi i\zeta},\qquad w=e^{-2\pi i\eta},\qquad
 zw=t,\qquad t=e^{2\pi i\tau}.
\]
Writing $\rho=e^{-2\pi h}$ and $r=e^{-2\pi s}$ gives the exact bounds
\[
 0<r<\rho<1,\qquad |t|=r^2<\rho^2=e^{-4\pi h}.
\]
The sewn collar is the annulus
\[
 \frac{|t|}{\rho}<|z|<\rho,\qquad
 \operatorname{mod}=\frac{1}{2\pi}\log\frac{\rho^2}{|t|}
                   =\operatorname{Im}\tau-2h.
\]
On its middle circle the sewing is
$re^{i\theta}\mapsto re^{i(\phi-\theta)}$, where
$t=|t|e^{i\phi}$. More generally, after rescaling by positive
$0<r_+,r_-<\rho$ satisfying $r_+r_-=|t|$, put $u=z/r_+$ and $v=w/r_-$.
Then the same sewing is
\[
 uv=e^{i\phi},\qquad
 t\,\partial_z\otimes\partial_w
   =e^{i\phi}\,\partial_u\otimes\partial_v.
\]
Thus the toric boundary welding in Proposition~\ref{A1-welding-plumbing}
is exactly this plumbing operation when its full coordinate germs are
retained. The normal tensor is its first-order invariant.
\end{proposition}
\begin{proof}
The exponential maps are invariant under $A$, have nonzero derivatives,
and give biholomorphisms from the two half-plane quotients to the
punctured disc of radius $\rho$. Their signs follow directly from
$|e^{2\pi i\zeta}|=e^{-2\pi\operatorname{Im}\zeta}$.
On the identified points $\zeta-\eta=\tau$, so multiplication gives
$zw=e^{2\pi i\tau}$. The two cutting lines
$\operatorname{Im}\zeta=s$ and $\operatorname{Im}\eta=-s$ map to
the circles of radius $r$. Since $|t|=e^{-2\pi\operatorname{Im}\tau}
=e^{-4\pi s}$, these radii and the collar inequalities are exact.

Deleting the smaller discs and identifying their boundary circles by
$w=t/z$ gives the same analytic surface as identifying the overlapping
annuli $|t|/\rho<|z|<\rho$ and $|t|/\rho<|w|<\rho$.
Indeed the transition is holomorphic with nonzero derivative and
continues across the seam; the two half-annuli form this one annulus.
Its modulus is the displayed logarithmic ratio. The phase formula and
rescaled sewing equation follow by substitution. Finally,
$\partial_u=r_+\partial_z$ and
$\partial_v=r_-\partial_w$, proving the tensor identity.
Under any further changes $z'=bz+O(z^2)$, $w'=cw+O(w^2)$,
the first-order smoothing scalar is multiplied by $bc$ while the
tangent basis is multiplied by $(bc)^{-1}$. Consequently the tensor
is independent of those coordinate descriptions. No assertion about
the higher coordinate jets follows from this first-order invariance.
\end{proof}

\paragraph{What was checked in the original.}
The source's unambiguous classical construction is $zw=s$ in
\S4.8.1, equation~(4.8.1), followed by cutting at radius $\sqrt{|s|}$
and sewing by $w=s/z$ in equation~(4.8.2). Its \S6.1.1, p.~67,
also writes the annular bundle explicitly as
$p(z,w)=zw$ on two punctured unit discs.
The printed p.~50 image was inspected: its upper/lower half-plane
exponential signs and its displayed angular expression are inconsistent
with its own disc and sewing conventions. The signs and angle in the
proposition are obtained by the direct calculation above; they are
not silently copied from those draft misprints. The exact
$e^{-4\pi h}$ bound follows from the exponential relation; the printed
$e^{-4h}$ bound is merely weaker.

For comparison, the annular-cover argument in \S4.1(III), p.~42,
equation~(4.1.1), places its cyclic cover between strips of widths
$\operatorname{Im}\tau-2h$ and $\operatorname{Im}\tau+2h$.
Whenever those stated strip inclusions hold, monotonicity of annulus
modulus and $\ell=\pi/\operatorname{mod}$ for the core of a complete
hyperbolic annulus give
\[
 \frac{\pi}{\operatorname{Im}\tau+2h}
 \leq\ell\leq
 \frac{\pi}{\operatorname{Im}\tau-2h}.
\]
In particular, for this fixed strip bound,
$\ell\log(1/|t|)\longrightarrow2\pi^2$ as
$\operatorname{Im}\tau\longrightarrow\infty$, the source's
Proposition~4.5.1, equation~(4.5.1), pp.~50--51.
The coordinate identification itself uses only the explicit local
construction proved above.

\paragraph{Exact source status and scope.}
This directly opened joint original discharges the requested
identification with Earle--Marden's plumbing construction.
Its global statements, including Theorem~4.6.1, p.~51, have not been
used as proof inputs here. They are distinct from the directly verified
local operation, and the draft itself warns in \S4.8.3, p.~58, that
arbitrary local coordinates do not give global plumbing coordinates.
The separate published chapter \emph{Holomorphic plumbing coordinates},
\emph{Contemporary Mathematics}~575 (2012), pp.~41--52, remains
unopened in this investigation. The draft is neither identified with
that chapter nor assigned its publication status.

